%% ============================================================
%%  Lecture 8 -- Tuesday, 23 June 2026
%%  Subfile of main.tex: compiles standalone OR as part of the book.
%%  Converted from math4061-current_notes/maths4061_lecture_08-2026_06_23/.
%% ============================================================
\documentclass[main.tex]{subfiles}

\begin{document}

\beginlecture{8}{Tuesday, 23 June 2026}%
  {Derivatives, Taylor's Theorem, and Analyticity}

\epigraph{And what are these fluxions? The velocities of evanescent increments.}
{George Berkeley, The Analyst (1734)}

\begin{lecturecontents}{Contents of this lecture}
  \item \hyperlink{8.1}{8.1\quad Derivatives as first-order approximations}
  \item \hyperlink{8.2}{8.2\quad Rules for derivatives}
  \item \hyperlink{8.3}{8.3\quad Oscillation examples}
  \item \hyperlink{8.4}{8.4\quad Extrema and the Mean Value Theorem}
  \item \hyperlink{8.5}{8.5\quad Darboux's theorem}
  \item \hyperlink{8.6}{8.6\quad Taylor polynomials and Taylor's theorem}
  \item \hyperlink{8.7}{8.7\quad Real and complex analyticity}
\end{lecturecontents}

\bigskip
\noindent\textbf{Overview.} This lecture begins the differential part of the course. The
derivative is first defined by the difference quotient, then immediately reinterpreted as the
best first-order linear approximation. That viewpoint makes the usual rules of differentiation
bookkeeping in powers of the increment, and it also explains the multivariable definition.
The lecture then turns the derivative into a global tool through Fermat's interior-extremum
theorem, Rolle's theorem, the Mean Value Theorem, monotonicity, and Darboux's theorem. It closes
with Taylor's theorem and the distinction between smooth and analytic functions, first over
$\mathbb R$ and then over $\mathbb C$. The companion text is Rudin, Chapter~5, with the
closest matches marked by the locators \R{5.1}--\R{5.16}.

\clearpage
% ============================================================
\sub{8.1}{Derivatives as first-order approximations}

\begin{Obj}{Definition}{8.1.1}{Derivative at a point}
Let $f:(a,b)\to\mathbb R$ and let $x_0\in(a,b)$. The function $f$ is
\emph{differentiable at $x_0$} (\R{5.1}) if the limit
\[
  \lim_{x\to x_0}\frac{f(x)-f(x_0)}{x-x_0}
  =
  \lim_{h\to0}\frac{f(x_0+h)-f(x_0)}{h}
\]
exists. This limit is the \emph{derivative} of $f$ at $x_0$, denoted $f'(x_0)$.
\end{Obj}

\begin{Fig}{8.1.2}{Best linear approximation}{Near $x_0$, differentiability says the
graph of $f$ is shadowed by an affine line whose error is smaller than the displacement.}
\begin{tikzpicture}[x=1cm,y=1cm]
  \def\xz{2.2}
  \def\xh{3.8}
  \pgfmathsetmacro{\fz}{0.18*(\xz-1)*(\xz-1)+1.1}
  \pgfmathsetmacro{\fh}{0.18*(\xh-1)*(\xh-1)+1.1}
  \pgfmathsetmacro{\slope}{0.36*(\xz-1)}
  \pgfmathsetmacro{\lin}{\fz+\slope*(\xh-\xz)}
  \pgfmathsetmacro{\lleft}{\fz+\slope*(0.3-\xz)}
  \draw[->] (-0.3,0) -- (6.0,0) node[right] {$x$};
  \draw[->] (0,-0.4) -- (0,4.4) node[above] {$y$};
  \draw[thick,figTerm,domain=0.35:5.25,smooth,variable=\x]
        plot ({\x},{0.18*(\x-1)*(\x-1)+1.1});
  \draw[thick,figLim,dashed] (0.3,\lleft) -- (\xh,\lin);
  \node[figLim,sloped,below=2pt,font=\small] at (1.4,{\fz+\slope*(1.4-\xz)})
        {$f(x_0)+ch$};
  \draw[dotted] (\xz,0) -- (\xz,\fz);
  \draw[dotted] (\xh,0) -- (\xh,\fh);
  \node[below] at (\xz,-0.03) {$x_0$};
  \node[below] at (\xh,-0.03) {$x_0+h$};
  \draw[<->,figLim] (\xz,-0.42) -- (\xh,-0.42) node[midway,below=5pt] {$h$};
  \pgfmathsetmacro{\xgap}{\xh+0.15}
  \draw[<->,thick,figBad] (\xgap,\lin) -- (\xgap,\fh)
        node[midway,right=6pt] {$o(h)$};
  \filldraw[figTerm] (\xz,\fz) circle (1.5pt);
  \filldraw[figTerm] (\xh,\fh) circle (1.5pt);
  \filldraw[figLim] (\xh,\lin) circle (1.5pt);
\end{tikzpicture}
\end{Fig}

\begin{Obj}{Proposition}{8.1.3}{Differentiability as a first-order expansion}
For $f:(a,b)\to\mathbb R$ and $x_0\in(a,b)$, the following are equivalent:
\[
  f\text{ is differentiable at }x_0
  \qquad\Longleftrightarrow\qquad
  f(x_0+h)=f(x_0)+ch+o(h)
\]
for some $c\in\mathbb R$, where $o(h)$ means an error $E(h)$ satisfying
\[
  \frac{E(h)}{h}\longrightarrow0\qquad(h\to0).
\]
In that case $c=f'(x_0)$.
This is Rudin's derivative definition \R{5.1} written in the increment variable.
\end{Obj}

\begin{Prf}{8.1.3}{Proof by Rearranging the Difference Quotient.}{[Direct]\quad increment form of \R{5.1}.}
\item Suppose $f$ is differentiable at $x_0$. Define
  \[
    E(h)=f(x_0+h)-f(x_0)-f'(x_0)h .
  \]
  Then the desired expansion holds with $c=f'(x_0)$, and for $h\ne0$,
  \[
    \frac{E(h)}{h}
    =
    \frac{f(x_0+h)-f(x_0)}{h}-f'(x_0)
    \longrightarrow0 .
  \]
\item Conversely, if $f(x_0+h)=f(x_0)+ch+E(h)$ and $E(h)/h\to0$, then
  \[
    \frac{f(x_0+h)-f(x_0)}{h}
    =
    c+\frac{E(h)}{h}
    \longrightarrow c .
  \]
  Hence the difference quotient converges and $f'(x_0)=c$.
\end{Prf}

\begin{Obj}{Definition}{8.1.4}{Differentiability from $\mathbb R^n$ to $\mathbb R^m$}
Let $F:\mathbb R^n\to\mathbb R^m$ and let $x_0\in\mathbb R^n$. The function $F$ is
\emph{differentiable at $x_0$} if there exists a linear map
$A:\mathbb R^n\to\mathbb R^m$ such that
\[
  F(x_0+h)=F(x_0)+A(h)+E(h),
  \qquad
  \frac{\|E(h)\|}{\|h\|}\longrightarrow0
  \quad(h\to0).
\]
The linear map $A$ is unique; in the standard bases it is represented by the Jacobian matrix.
Compare Rudin's one-parameter vector-valued derivative in \R{5.16}; the present definition
replaces one real increment by a vector increment.
\end{Obj}

\begin{ObjL}{Remark}{8.1.5}{Why the linear map replaces division}
In one real variable, the quotient
\[
  \frac{f(x_0+h)-f(x_0)}{h}
\]
is available because nonzero real increments have inverses. A vector increment
$h\in\mathbb R^n$ has no comparable division operation, so the higher-dimensional definition
subtracts a candidate linear map and asks that the remaining vector error be smaller than
$\|h\|$. Complex differentiability is a special two-dimensional exception: identifying
$\mathbb R^2$ with $\mathbb C$ restores division by a nonzero increment, but imposes a much
stronger condition than ordinary real differentiability. Compare the vector-valued
one-parameter derivative in \R{5.16}.
\end{ObjL}

% ============================================================
\sub{8.2}{Rules for derivatives}

\begin{Obj}{Theorem}{8.2.1}{Differentiability implies continuity}
If $f$ is differentiable at $x_0$, then $f$ is continuous at $x_0$ (\R{5.2}).
\end{Obj}

\begin{Prf}{8.2.1}{Proof from the First-Order Expansion.}{[Direct]\quad cf.\ \R{5.2}.}
\item By \xref{8.1.3},
  \[
    f(x_0+h)=f(x_0)+f'(x_0)h+E(h),
    \qquad \frac{E(h)}{h}\to0 .
  \]
\item Since $E(h)=h(E(h)/h)$, we have $E(h)\to0$. Therefore
  \[
    f(x_0+h)-f(x_0)=f'(x_0)h+E(h)\longrightarrow0,
  \]
  which is continuity at $x_0$.
\end{Prf}

\begin{Obj}{Corollary}{8.2.2}{Discontinuity obstructs differentiability}
If $f$ is not continuous at $x_0$, then $f$ is not differentiable at $x_0$.
This is the contrapositive of Rudin's differentiability-implies-continuity theorem \R{5.2}.
\end{Obj}

\begin{Prf}{8.2.2}{Proof by Contrapositive.}{[A9]}
\item The contrapositive of \xref{8.2.1} says precisely that differentiability at $x_0$
  forces continuity at $x_0$.
\item Therefore a failure of continuity at $x_0$ rules out differentiability there.
\end{Prf}

\begin{Obj}{Proposition}{8.2.3}{Algebra rules for derivatives}
Suppose $f$ and $g$ are differentiable at $x_0$. Then
\[
  (f\pm g)'(x_0)=f'(x_0)\pm g'(x_0),
  \qquad
  (cf)'(x_0)=c\,f'(x_0),
\]
\[
  (fg)'(x_0)=f'(x_0)g(x_0)+f(x_0)g'(x_0).
\]
If $g(x_0)\ne0$, then $f/g$ is differentiable at $x_0$ and
\[
  \left(\frac{f}{g}\right)'(x_0)
  =
  \frac{f'(x_0)g(x_0)-f(x_0)g'(x_0)}{g(x_0)^2}.
\]
These are Rudin's algebra rules for derivatives (\R{5.3}).
\end{Obj}

\begin{Prf}{8.2.3}{Proof by First-Order Bookkeeping.}{[Direct]\quad cf.\ \R{5.3}.}
\item Write
  \[
    f(x_0+h)=f(x_0)+f'(x_0)h+E_f(h),
    \qquad
    g(x_0+h)=g(x_0)+g'(x_0)h+E_g(h),
  \]
  with $E_f(h)/h\to0$ and $E_g(h)/h\to0$.
\item Adding the expansions gives
  \[
    (f+g)(x_0+h)
    =
    (f+g)(x_0)+\bigl(f'(x_0)+g'(x_0)\bigr)h+\bigl(E_f(h)+E_g(h)\bigr),
  \]
  and the last error is $o(h)$ because
  \[
    \frac{E_f(h)+E_g(h)}{h}
    =
    \frac{E_f(h)}{h}+\frac{E_g(h)}{h}\longrightarrow0 .
  \]
  The difference and scalar-multiple rules are the same calculation.
\item Multiplying the expansions and collecting the constant and first-order terms gives
  \[
  \begin{aligned}
    (fg)(x_0+h)
    &=f(x_0)g(x_0)
      +f'(x_0)g(x_0)h
      +f(x_0)g'(x_0)h\\
    &\quad
      +E_f(h)\bigl(g(x_0)+g'(x_0)h+E_g(h)\bigr)\\
    &\quad
      +E_g(h)\bigl(f(x_0)+f'(x_0)h\bigr)
      +f'(x_0)g'(x_0)h^2 .
  \end{aligned}
  \]
  After division by $h$, the last line tends to $0$ because the error ratios tend to $0$
  and the bracketed factors remain bounded. Hence the coefficient of $h$ is the derivative
  of $fg$.
\item If $g(x_0)\ne0$, then by \xref{8.2.1} the values $g(x_0+h)$ stay nonzero for all
  sufficiently small $h$. Thus
  \[
  \begin{aligned}
    \left(\frac1g\right)'(x_0)
    &=
    \lim_{h\to0}\frac{1}{h}
    \left(\frac1{g(x_0+h)}-\frac1{g(x_0)}\right)\\
    &=
    -\lim_{h\to0}
      \frac{g(x_0+h)-g(x_0)}{h}\cdot
      \frac1{g(x_0+h)g(x_0)}
     =
     -\frac{g'(x_0)}{g(x_0)^2}.
  \end{aligned}
  \]
  Applying the product rule to $f\cdot(1/g)$ gives the quotient rule.
\end{Prf}

\begin{Fig}{8.2.4}{The product rule}{The two first-order strips contribute to the
derivative of the area; the corner is second-order.}
\begin{tikzpicture}[x=1cm,y=1cm]
  \def\W{3.6}\def\dW{1.0}\def\H{2.2}\def\dH{0.75}
  \fill[figTerm!10] (0,0) rectangle (\W,\H);
  \draw[thick] (0,0) rectangle (\W,\H);
  \fill[figLim!16] (0,\H) rectangle (\W,\H+\dH);
  \fill[figLim!16] (\W,0) rectangle (\W+\dW,\H);
  \fill[figBad!10] (\W,\H) rectangle (\W+\dW,\H+\dH);
  \draw (0,\H) rectangle (\W,\H+\dH);
  \draw (\W,0) rectangle (\W+\dW,\H);
  \draw (\W,\H) rectangle (\W+\dW,\H+\dH);
  \node at (\W/2,\H/2) {$f(x_0)g(x_0)$};
  \node[font=\small] at (\W/2,\H+\dH/2) {$f(x_0)g'(x_0)h$};
  \node[font=\small,rotate=90] at (\W+\dW/2,\H/2) {$g(x_0)f'(x_0)h$};
  \node[font=\scriptsize] at (\W+\dW/2,\H+\dH/2) {$O(h^2)$};
  \node[below] at (\W/2,-0.08) {$f(x_0)$};
  \node[below] at (\W+\dW/2,-0.08) {$f'(x_0)h$};
  \node[left] at (-0.08,\H/2) {$g(x_0)$};
  \node[left] at (-0.08,\H+\dH/2) {$g'(x_0)h$};
\end{tikzpicture}
\end{Fig}

\begin{Obj}{Proposition}{8.2.5}{Chain rule}
Suppose $f$ is differentiable at $x_0$ and $g$ is differentiable at $f(x_0)$. Then
$g\circ f$ is differentiable at $x_0$ and
\[
  (g\circ f)'(x_0)=g'\bigl(f(x_0)\bigr)f'(x_0).
\]
This is the chain rule (\R{5.5}).
\end{Obj}

\begin{Prf}{8.2.5}{Proof by Composing First-Order Errors.}{[Direct]\quad cf.\ \R{5.5}.}
\item Let
  \[
    k(h)=f(x_0+h)-f(x_0).
  \]
  Differentiability of $f$ gives $k(h)=f'(x_0)h+E_1(h)$, with $E_1(h)/h\to0$; in particular
  $k(h)\to0$.
\item Differentiability of $g$ at $f(x_0)$ means there is an error $E_2(u)$ such that
  \[
    g(f(x_0)+u)=g(f(x_0))+g'(f(x_0))u+E_2(u),
    \qquad \frac{E_2(u)}{u}\to0 .
  \]
  To avoid dividing by a possibly zero displacement, define
  \[
    H(u)=
    \begin{cases}
      E_2(u)/u, & u\ne0,\\
      0, & u=0.
    \end{cases}
  \]
  Then $H(u)\to0$ as $u\to0$ and $E_2(u)=H(u)u$ for every $u$.
\item Substituting $u=k(h)$ gives
  \[
  \begin{aligned}
    g(f(x_0+h))
    &=
    g(f(x_0))+g'(f(x_0))k(h)+H(k(h))k(h)\\
    &=
    g(f(x_0))+g'(f(x_0))f'(x_0)h
      +g'(f(x_0))E_1(h)+H(k(h))k(h).
  \end{aligned}
  \]
\item The remainder divided by $h$ tends to $0$:
  \[
    g'(f(x_0))\frac{E_1(h)}{h}
    +
    H(k(h))\frac{k(h)}{h}
    \longrightarrow
    0+0\cdot f'(x_0)=0.
  \]
  Therefore the coefficient of $h$ is the derivative of $g\circ f$ at $x_0$.
\end{Prf}

\begin{Obj}{Proposition}{8.2.6}{Power rule}
For every integer $n\ge0$,
\[
  \frac{d}{dx}x^n=nx^{n-1},
\]
where the case $n=0$ is read separately as the derivative of the constant function $1$,
namely $0$. \R{5.4}
\end{Obj}

\begin{Prf}{8.2.6}{Proof by Induction Using the Product Rule.}{[A10]}
\item For $n=0$, the function $x^0=1$ is constant, so its derivative is $0$.
\item For $n=1$, the difference quotient of $x$ is $1$.
\item Assume $(x^n)'=nx^{n-1}$. Since $x^{n+1}=x\cdot x^n$, the product rule gives
  \[
    (x^{n+1})'
    =
    (x)'x^n+x(x^n)'
    =
    x^n+nx^n
    =
    (n+1)x^n .
  \]
  The formula follows by induction.
\end{Prf}

% ============================================================
\sub{8.3}{Oscillation examples}

\begin{Fig}{8.3.1}{Three oscillatory functions}{Multiplying by powers of $x$ progressively
tames the oscillation of $\sin(1/x)$ near $0$.}
\begin{tikzpicture}[x=1cm,y=1cm]
  \begin{scope}[shift={(0,0)},xscale=1.7,yscale=0.9]
    \draw[->] (-0.8,0) -- (0.8,0);
    \draw[->] (0,-1.2) -- (0,1.25);
    \draw[figTerm,domain=0.045:0.7,samples=260,smooth] plot (\x,{sin(deg(1/\x))});
    \draw[figTerm,domain=-0.7:-0.045,samples=260,smooth] plot (\x,{sin(deg(1/\x))});
    \node[font=\small] at (0,1.45) {$\sin(1/x)$};
    \node[font=\scriptsize,align=center] at (0,-1.45) {not continuous};
  \end{scope}
  \begin{scope}[shift={(3.3,0)},xscale=1.7,yscale=1.8]
    \draw[->] (-0.8,0) -- (0.8,0);
    \draw[->] (0,-0.75) -- (0,0.82);
    \draw[figLim,dashed] (-0.7,-0.7) -- (0.7,0.7);
    \draw[figLim,dashed] (-0.7,0.7) -- (0.7,-0.7);
    \draw[figTerm,domain=0.045:0.7,samples=260,smooth] plot (\x,{\x*sin(deg(1/\x))});
    \draw[figTerm,domain=-0.7:-0.045,samples=260,smooth] plot (\x,{\x*sin(deg(1/\x))});
    \node[font=\small] at (0,0.98) {$x\sin(1/x)$};
    \node[font=\scriptsize,align=center] at (0,-0.95) {continuous,\\not differentiable};
  \end{scope}
  \begin{scope}[shift={(6.6,0)},xscale=1.7,yscale=2.6]
    \draw[->] (-0.8,0) -- (0.8,0);
    \draw[->] (0,-0.55) -- (0,0.62);
    \draw[figLim,dashed,domain=-0.7:0.7,samples=80] plot (\x,{\x*\x});
    \draw[figLim,dashed,domain=-0.7:0.7,samples=80] plot (\x,{-\x*\x});
    \draw[figTerm,domain=0.045:0.7,samples=260,smooth] plot (\x,{\x*\x*sin(deg(1/\x))});
    \draw[figTerm,domain=-0.7:-0.045,samples=260,smooth] plot (\x,{\x*\x*sin(deg(1/\x))});
    \node[font=\small] at (0,0.76) {$x^2\sin(1/x)$};
    \node[font=\scriptsize,align=center] at (0,-0.72) {differentiable,\\$f'$ not continuous};
  \end{scope}
\end{tikzpicture}
\end{Fig}

\begin{ObjL}{Example}{8.3.2}{$\sin(1/x)$ is not continuous at $0$}
Define $f(0)=0$ and $f(x)=\sin(1/x)$ for $x\ne0$. Along sequences approaching $0$, the values
can approach different limits; for example one may choose points where $\sin(1/x)=1$ and
others where $\sin(1/x)=-1$. Thus $f$ is not continuous at $0$, and by \xref{8.2.2} it is not
differentiable there (cf.\ \R{5.6}).
\end{ObjL}

\begin{ObjL}{Example}{8.3.3}{$x\sin(1/x)$ is continuous but not differentiable at $0$}
Let $f(0)=0$ and $f(x)=x\sin(1/x)$ for $x\ne0$. Since
\[
  |x\sin(1/x)|\le |x|,
\]
the squeeze theorem gives $f(x)\to0=f(0)$. But the difference quotient at $0$ is
\[
  \frac{f(h)-f(0)}{h}=\sin(1/h),
\]
which has no limit. Hence $f$ is continuous at $0$ but not differentiable there (\R{5.6}).
\end{ObjL}

\begin{ObjL}{Example}{8.3.4}{$x^2\sin(1/x)$ is differentiable but not $C^1$ at $0$}
Let $f(0)=0$ and $f(x)=x^2\sin(1/x)$ for $x\ne0$. Then
\[
  \frac{f(h)-f(0)}{h}=h\sin(1/h)\longrightarrow0,
\]
so $f'(0)=0$. For $x\ne0$, the product and chain rules give
\[
  f'(x)=2x\sin(1/x)-\cos(1/x).
\]
The first term tends to $0$, but $\cos(1/x)$ has no limit as $x\to0$. Thus $f$ is
differentiable at $0$, yet $f'$ is not continuous at $0$ (\R{5.6}).
\end{ObjL}

\begin{ObjL}{Remark}{8.3.5}{The hierarchy already separates}
The three examples show that differentiability is stronger than continuity, and continuous
differentiability is stronger than differentiability. Each extra factor of $x$ narrows the
envelope around the oscillation and buys one more degree of regularity at the origin
(cf.\ \R{5.6}).
\end{ObjL}

% ============================================================
\sub{8.4}{Extrema and the Mean Value Theorem}

\begin{Obj}{Definition}{8.4.1}{Local extrema}
Let $f:E\to\mathbb R$ and let $p\in E$. The function $f$ has a \emph{local maximum} at
$p$ if there exists $\delta>0$ such that
\[
  q\in E,\quad |q-p|<\delta
  \quad\Longrightarrow\quad
  f(q)\le f(p).
\]
A \emph{local minimum} is defined by reversing the inequality (\R{5.7}).
\end{Obj}

\begin{Fig}{8.4.2}{A local maximum}{On a sufficiently small window around $p$, the value
$f(p)$ dominates nearby values of the function.}
\begin{tikzpicture}[x=1cm,y=1cm]
  \draw[->] (-3,0) -- (3,0);
  \draw[thick,figTerm,domain=-2.7:2.7,samples=180,smooth]
        plot (\x,{1.65*exp(-0.85*\x*\x)});
  \filldraw[figTerm] (0,1.65) circle (1.5pt);
  \draw[dashed,figLim] (-1.35,1.65) -- (1.35,1.65);
  \draw[dotted] (0,1.65) -- (0,0);
  \foreach \x in {-1.35,0,1.35} \draw (\x,0.07) -- (\x,-0.07);
  \node[below] at (-1.35,-0.08) {$p-\delta$};
  \node[below] at (0,-0.08) {$p$};
  \node[below] at (1.35,-0.08) {$p+\delta$};
  \draw[<->,figLim]
        (-1.35,-0.55) -- (1.35,-0.55)
        node[midway,below=5pt] {$f(q)\le f(p)$ nearby};
\end{tikzpicture}
\end{Fig}

\begin{Obj}{Theorem}{8.4.3}{Interior extremum theorem}
Let $f:(a,b)\to\mathbb R$ be differentiable. If $x_0\in(a,b)$ is a local maximum or a
local minimum, then
\[
  f'(x_0)=0.
\]
This is Fermat's interior-extremum theorem (\R{5.8}).
\end{Obj}

\begin{Prf}{8.4.3}{Proof by One-Sided Difference Quotients.}{[A12]\quad cf.\ \R{5.8}.}
\item Assume $x_0$ is a local maximum. For $x$ sufficiently close to $x_0$,
  \[
    f(x)-f(x_0)\le0 .
  \]
\item If $x<x_0$, then $x-x_0<0$, so
  \[
    \frac{f(x)-f(x_0)}{x-x_0}\ge0.
  \]
  Passing to the left-hand limit gives $f'(x_0)\ge0$.
\item If $x>x_0$, then $x-x_0>0$, so
  \[
    \frac{f(x)-f(x_0)}{x-x_0}\le0.
  \]
  Passing to the right-hand limit gives $f'(x_0)\le0$.
\item The two one-sided limits are equal because $f'(x_0)$ exists. Therefore
  \[
    0\le f'(x_0)\le0,
  \]
  and $f'(x_0)=0$. The minimum case is identical with the inequalities reversed.
\end{Prf}

\begin{ObjL}{Remark}{8.4.4}{Why the hypotheses matter}
The point must be interior because the proof uses both sides. At an endpoint only one one-sided
inequality is available. The converse also fails: $f(x)=x^3$ has $f'(0)=0$, but $0$ is not a
local extremum (cf.\ \R{5.8}).
\end{ObjL}

\begin{Obj}{Proposition}{8.4.5}{Rolle's theorem}
Let $f:[a,b]\to\mathbb R$ be continuous on $[a,b]$ and differentiable on $(a,b)$. If
$f(a)=f(b)$, then there exists $c\in(a,b)$ such that
\[
  f'(c)=0 .
\]
This is the Rolle special case behind Rudin's generalized mean-value theorem (\R{5.9}).
\end{Obj}

\begin{Fig}{8.4.6}{Rolle's theorem}{Equal endpoint heights force a horizontal tangent at
some interior point, unless the function is constant.}
\begin{tikzpicture}[x=1cm,y=1cm]
  \draw[->] (-2.7,0) -- (2.7,0);
  \draw[dashed,figLim] (-2,0.6) -- (2,0.6);
  \draw[thick,figTerm,domain=-2:2,samples=120,smooth] plot (\x,{1.8-0.3*\x*\x});
  \filldraw[figTerm] (-2,0.6) circle (1.4pt);
  \filldraw[figTerm] (2,0.6) circle (1.4pt);
  \filldraw[figLim] (0,1.8) circle (1.4pt);
  \draw[dashed,figLim] (-1.05,1.8) -- (1.05,1.8);
  \draw[dotted] (-2,0.6) -- (-2,0);
  \draw[dotted] (2,0.6) -- (2,0);
  \draw[dotted] (0,1.8) -- (0,0);
  \foreach \x in {-2,0,2} \draw (\x,0.07) -- (\x,-0.07);
  \node[below] at (-2,-0.08) {$a$};
  \node[below] at (0,-0.08) {$c$};
  \node[below] at (2,-0.08) {$b$};
  \node[above,figLim] at (0,1.86) {$f'(c)=0$};
\end{tikzpicture}
\end{Fig}

\begin{Prf}{8.4.5}{Proof by the Extreme-Value Theorem and Interior Extrema.}{[A6$+$A12]\quad Rolle step behind \R{5.9}.}
\item Since $[a,b]$ is compact and $f$ is continuous, $f$ attains a maximum and a minimum on
  $[a,b]$ by the extreme-value theorem from Lecture~7.
\item If either extreme is attained at an interior point $c\in(a,b)$, then \xref{8.4.3}
  gives $f'(c)=0$.
\item If no extreme is attained in the interior, then both the maximum and minimum occur at
  endpoints. Since $f(a)=f(b)$, the maximum and minimum values coincide. Hence $f$ is
  constant on $[a,b]$, and any $c\in(a,b)$ satisfies $f'(c)=0$.
\end{Prf}

\begin{Obj}{Corollary}{8.4.7}{Mean Value Theorem}
Let $f:[a,b]\to\mathbb R$ be continuous on $[a,b]$ and differentiable on $(a,b)$. Then
there exists $c\in(a,b)$ such that
\[
  f'(c)=\frac{f(b)-f(a)}{b-a}.
\]
This is the Mean Value Theorem (\R{5.10}).
\end{Obj}

\begin{Fig}{8.4.8}{Mean Value Theorem}{Some interior tangent is parallel to the chord
joining the endpoints.}
\begin{tikzpicture}[x=1cm,y=1cm]
  \draw[->] (-2.8,0) -- (3.0,0);
  \draw[thick,figTerm,domain=-2.1:2.4,samples=150,smooth]
        plot (\x,{1.6-0.18*\x*\x+0.35*\x});
  \draw[figLim,thick] (-2.35,{0.35*(-2.35)+0.88}) -- (2.75,{0.35*2.75+0.88});
  \draw[figLim,dashed] (-1.2,{0.35*(-1.2)+1.6}) -- (1.25,{0.35*1.25+1.6});
  \draw (-2,0.18) circle (2pt);
  \draw (2,1.58) circle (2pt);
  \filldraw[figLim] (0,1.6) circle (1.4pt);
  \draw[dotted] (-2,0.18) -- (-2,0);
  \draw[dotted] (2,1.58) -- (2,0);
  \draw[dotted] (0,1.6) -- (0,0);
  \foreach \x in {-2,0,2} \draw (\x,0.07) -- (\x,-0.07);
  \node[below] at (-2,-0.08) {$a$};
  \node[below] at (0,-0.08) {$c$};
  \node[below] at (2,-0.08) {$b$};
\end{tikzpicture}
\end{Fig}

\begin{Prf}{8.4.7}{Proof by Tilting the Chord Flat.}{[Constr]\quad cf.\ \R{5.10}.}
\item Let $L$ be the secant line
  \[
    L(x)=f(a)+\frac{f(b)-f(a)}{b-a}(x-a),
  \]
  and define $h=f-L$.
\item The function $h$ is continuous on $[a,b]$ and differentiable on $(a,b)$. Also
  \[
    h(a)=0,\qquad h(b)=0.
  \]
\item By Rolle's theorem, there exists $c\in(a,b)$ with $h'(c)=0$. Since
  \[
    h'(x)=f'(x)-\frac{f(b)-f(a)}{b-a},
  \]
  the conclusion follows.
\end{Prf}

\begin{Obj}{Corollary}{8.4.9}{Monotonicity test}
Let $f$ be differentiable on $(a,b)$.
\begin{enumerate}
  \item If $f'(x)\ge0$ for all $x\in(a,b)$, then $f$ is nondecreasing.
  \item If $f'(x)=0$ for all $x\in(a,b)$, then $f$ is constant.
  \item If $f'(x)\le0$ for all $x\in(a,b)$, then $f$ is nonincreasing.
\end{enumerate}
These are the standard monotonicity consequences of the Mean Value Theorem (\R{5.11}).
\end{Obj}

\begin{Prf}{8.4.9}{Proof by Applying the Mean Value Theorem on Subintervals.}{[Direct]\quad cf.\ \R{5.11}.}
\item Fix $x_1<x_2$ in $(a,b)$. By \xref{8.4.7}, there exists $\xi\in(x_1,x_2)$ such that
  \[
    f(x_2)-f(x_1)=f'(\xi)(x_2-x_1).
  \]
\item Since $x_2-x_1>0$, the sign of $f(x_2)-f(x_1)$ is the sign of $f'(\xi)$. The three
  conclusions follow immediately from the three sign hypotheses.
\end{Prf}

% ============================================================
\sub{8.5}{Darboux's theorem}

\begin{Obj}{Proposition}{8.5.1}{Darboux's theorem}
Let $f$ be differentiable on an interval containing $a<b$. If $\lambda$ lies strictly
between $f'(a)$ and $f'(b)$, then there exists $c\in(a,b)$ such that
\[
  f'(c)=\lambda .
\]
Thus derivatives have the intermediate value property, even though derivatives need not be
continuous (\R{5.12}).
\end{Obj}

\begin{Prf}{8.5.1}{Proof by Tilting and Taking an Interior Extremum.}{[A12$+$Constr]\quad cf.\ \R{5.12}.}
\item Assume $f'(a)<\lambda<f'(b)$; the reverse case is obtained by changing signs. Define
  \[
    g(x)=f(x)-\lambda x.
  \]
  Then $g'(a)<0$ and $g'(b)>0$.
\item Since $g$ is continuous on $[a,b]$, it attains a minimum at some point $c\in[a,b]$.
\item The condition $g'(a)<0$ means that just to the right of $a$ the values of $g$ are
  smaller than $g(a)$, so the minimum is not attained at $a$. Similarly, $g'(b)>0$ implies
  that just to the left of $b$ the values of $g$ are smaller than $g(b)$, so the minimum is
  not attained at $b$.
\item Hence the minimum is attained at an interior point $c\in(a,b)$. By \xref{8.4.3},
  $g'(c)=0$, and therefore
  \[
    f'(c)-\lambda=0.
  \]
\end{Prf}

\begin{ObjL}{Remark}{8.5.2}{What Darboux rules out}
Darboux's theorem forbids derivatives from having jump discontinuities, because a jump would
skip the intermediate values. The theorem does not say that derivatives are continuous:
\xref{8.3.4} gives a derivative with an oscillatory discontinuity at $0$ (cf.\ \R{5.12}).
\end{ObjL}

% ============================================================
\sub{8.6}{Taylor polynomials and Taylor's theorem}

\begin{Obj}{Definition}{8.6.1}{Taylor polynomial}
Let $n\ge1$, and suppose $f$ has derivatives through order $n-1$ at $\alpha$. The
\emph{Taylor polynomial of degree at most $n-1$ centered at $\alpha$} is
\[
  P_{n-1}(x)=\sum_{k=0}^{n-1}\frac{f^{(k)}(\alpha)}{k!}(x-\alpha)^k,
  \qquad f^{(0)}=f.
\]
Rudin introduces higher derivatives in \R{5.14} and this polynomial in \R{5.15}.
\end{Obj}

\begin{Obj}{Proposition}{8.6.2}{Matching and uniqueness}
The Taylor polynomial $P_{n-1}$ is the unique polynomial of degree at most $n-1$ satisfying
\[
  P_{n-1}^{(k)}(\alpha)=f^{(k)}(\alpha),
  \qquad k=0,1,\dots,n-1.
\]
This is the matching property built into Rudin's Taylor polynomial in \R{5.15}.
\end{Obj}

\begin{Prf}{8.6.2}{Proof in the Shifted Polynomial Basis.}{[Direct]}
\item Write a general degree-at-most $n-1$ polynomial in the shifted basis:
  \[
    Q(x)=\sum_{k=0}^{n-1}c_k(x-\alpha)^k .
  \]
\item Differentiating $k$ times and evaluating at $\alpha$ leaves only the $k$-th shifted
  monomial:
  \[
    Q^{(k)}(\alpha)=k!\,c_k.
  \]
\item The matching conditions therefore force
  \[
    c_k=\frac{f^{(k)}(\alpha)}{k!},
  \]
  exactly the coefficients of $P_{n-1}$. Hence the matching polynomial exists and is unique.
\end{Prf}

\begin{Obj}{Theorem}{8.6.3}{Taylor's theorem with Lagrange remainder}
Let $n\ge1$. Suppose $f:[a,b]\to\mathbb R$ has $f^{(n-1)}$ continuous on $[a,b]$ and
$f^{(n)}$ exists on $(a,b)$. If $\alpha,\beta\in[a,b]$ are distinct and $P_{n-1}$ is the
Taylor polynomial centered at $\alpha$, then there is a point $\xi$ strictly between
$\alpha$ and $\beta$ such that
\[
  f(\beta)
  =
  P_{n-1}(\beta)+
  \frac{f^{(n)}(\xi)}{n!}(\beta-\alpha)^n .
\]
This is Rudin's Taylor theorem with Lagrange remainder (\R{5.15}).
\end{Obj}

\begin{Prf}{8.6.3}{Proof by Iterated Rolle's Theorem.}{[A6$+$Constr]\quad cf.\ \R{5.15}.}
\item Assume $\alpha<\beta$; the other order is identical on the interval between them.
  Choose $M$ by
  \[
    M=\frac{f(\beta)-P_{n-1}(\beta)}{(\beta-\alpha)^n},
  \]
  so that $f(\beta)=P_{n-1}(\beta)+M(\beta-\alpha)^n$.
\item Define
  \[
    g(x)=f(x)-P_{n-1}(x)-M(x-\alpha)^n .
  \]
  By the matching property in \xref{8.6.2},
  \[
    g(\alpha)=g'(\alpha)=\cdots=g^{(n-1)}(\alpha)=0,
  \]
  and by the choice of $M$, $g(\beta)=0$.
\item Apply Rolle's theorem repeatedly. First, $g(\alpha)=g(\beta)=0$ gives a point
  $c_1\in(\alpha,\beta)$ with $g'(c_1)=0$. Then $g'(\alpha)=g'(c_1)=0$ gives a point
  $c_2\in(\alpha,c_1)$ with $g''(c_2)=0$. Continuing $n$ times gives a point
  $\xi\in(\alpha,\beta)$ with
  \[
    g^{(n)}(\xi)=0.
  \]
\item Since $P_{n-1}^{(n)}=0$ and $\frac{d^n}{dx^n}M(x-\alpha)^n=n!M$,
  \[
    0=g^{(n)}(\xi)=f^{(n)}(\xi)-n!M.
  \]
  Thus $M=f^{(n)}(\xi)/n!$, and substituting this into the definition of $M$ gives the
  stated formula.
\end{Prf}

\begin{ObjL}{Remark}{8.6.4}{The qualitative little-$o$ form}
The Lagrange formula gives the quantitative error. If, in addition, $f^{(n)}$ is locally
bounded near $\alpha$--for instance if $f^{(n)}$ is continuous near $\alpha$--then
\[
  f(t)-P_{n-1}(t)=O\bigl((t-\alpha)^n\bigr),
\]
and hence
\[
  \frac{f(t)-P_{n-1}(t)}{(t-\alpha)^{n-1}}\longrightarrow0 .
\]
The boundedness condition is doing real work here; the existence of $f^{(n)}$ alone does not
by itself provide local boundedness. This is the qualitative consequence of the Lagrange
remainder in \R{5.15}.
\end{ObjL}

% ============================================================
\sub{8.7}{Real and complex analyticity}

\begin{Obj}{Definition}{8.7.1}{Real analyticity at a point}
A smooth real function $f$ is \emph{real analytic at $\alpha$} if its Taylor series at
$\alpha$ converges to $f$ on some neighbourhood of $\alpha$:
\[
  f(t)=\sum_{k=0}^{\infty}\frac{f^{(k)}(\alpha)}{k!}(t-\alpha)^k
\]
for all $t$ sufficiently close to $\alpha$.
This goes beyond Rudin's finite Taylor theorem \R{5.15}: one must also control the limiting
behaviour of the remainders as the degree tends to infinity.
\end{Obj}

\begin{Obj}{Counterexample}{8.7.2}{A smooth function need not be analytic}
Define
\[
  f(x)=
  \begin{cases}
    e^{-1/x^2}, & x\ne0,\\
    0, & x=0.
  \end{cases}
\]
This function is $C^\infty$, and all derivatives at $0$ vanish. Its Taylor series at $0$ is
therefore identically $0$, but $f(x)>0$ for $x\ne0$. Hence smoothness does not imply real
analyticity. This is beyond Rudin's finite Taylor theorem \R{5.15}.
\end{Obj}
\fails{the claim that knowing all derivatives at one point determines a smooth real function.}

\begin{Obj}{Theorem}{8.7.3}{Weierstrass Function}
There exists a continuous function $W:\mathbb R\to\mathbb R$ that is differentiable at no
point. This is an advanced separation result beyond Rudin Chapter~5.
\end{Obj}

\PrfRef{8.7.3}{Quoted Advanced Result.}{This theorem is stated as part of the regularity
landscape; its proof is beyond the lecture.}

\begin{Fig}{8.7.4}{The real regularity tower}{For real functions the inclusions between
regularity classes are strict; each rung contains functions not lying on the next rung.}
\begin{tikzpicture}[x=1cm,y=1cm]
  \draw[thick,figTerm] (0,0) ellipse (4.7 and 2.9);
  \draw[thick,figTerm] (0,0) ellipse (3.65 and 2.25);
  \draw[thick,figTerm] (0,0) ellipse (2.55 and 1.55);
  \fill[figLim!12] (0,0) ellipse (1.45 and 0.85);
  \draw[thick,figLim] (0,0) ellipse (1.45 and 0.85);
  \node at (0,2.55) {$C^0$};
  \node at (0,1.95) {$C^1$};
  \node at (0,1.22) {$C^\infty$};
  \node[figLim] at (0,0.42) {$C^\omega$};
  \node[font=\small] at (0,-0.25) {$e^x,\ \sin x,\ \cos x$};
  \node[font=\small] at (0,-1.1) {$e^{-1/x^2}$};
  \node[font=\small] at (0,-1.82) {$C^k\setminus C^{k+1}$ examples};
  \node[font=\small] at (0,-2.55) {Weierstrass function};
\end{tikzpicture}
\end{Fig}

\begin{ObjL}{Remark}{8.7.5}{Derivative growth and Taylor convergence}
Taylor convergence is controlled by the growth of the coefficients
$f^{(r)}(\alpha)/r!$. If these coefficients grow too quickly, the formal Taylor series can
have radius of convergence $0$. A useful sufficient condition in the other direction is
uniform control of derivatives on a neighbourhood: if the Lagrange remainders satisfy a bound
forcing
\[
  \left|\frac{f^{(n)}(\xi)}{n!}(t-\alpha)^n\right|\longrightarrow0,
\]
then the Taylor polynomials converge back to $f$ near $\alpha$. Compare the finite
Lagrange-remainder control in \R{5.15}.
\end{ObjL}

\begin{Obj}{Theorem}{8.7.6}{Holomorphic functions are analytic}
Let $U\subseteq\mathbb C$ be open. If $f:U\to\mathbb C$ is complex differentiable at every
point of $U$, then $f$ is complex analytic: near each point of $U$, it is represented by a
convergent power series. This complex-analysis theorem is beyond Rudin Chapter~5.
\end{Obj}

\PrfRef{8.7.6}{Quoted Theorem from Complex Analysis.}{The result is recorded to contrast
real and complex differentiability; its proof belongs to complex analysis.}

\begin{ObjL}{Remark}{8.7.7}{The real tragedy and the complex collapse}
Over $\mathbb R$, the classes $C^1,C^2,\ldots,C^\infty,C^\omega$ do not collapse. Over
$\mathbb C$, the hypothesis is not ordinary real differentiability of a map
$\mathbb R^2\to\mathbb R^2$; it is complex differentiability on an open set, a much more rigid
condition. Under that holomorphic hypothesis, one derivative already forces local power-series
representation. This is the contrast with Rudin's finite Taylor theorem \R{5.15} and beyond the
real-variable scope of Chapter~5.
\end{ObjL}

% per-lecture Index of Objects -- standalone compile only; the book uses the master index.
\ifSubfilesClassLoaded{\clearpage\objindex}{}

\end{document}
